\documentclass[11pt,a4paper]{article}

\usepackage[utf8]{inputenc}
\usepackage[T1]{fontenc}
\usepackage{lmodern}
\usepackage[margin=2.5cm]{geometry}
\usepackage{hyperref}

\hypersetup{colorlinks=true, urlcolor=blue}

\title{Question \#2: Analytical Reasoning \\ Pose Estimation for Padel Hits}
\author{Alejandro Campayo}
\date{\today}

\begin{document}

\maketitle

\section*{Q1. Root causes of validation-to-new-club performance drop}

Strong performance on validation and a failure in a new club means that the difference
between \textbf{validation and training data} is smaller than the difference between
\textbf{training and real deployment}: there is a \textbf{domain shift}. What worries me
most is that the validation metrics do not show it, which points to \textbf{data leakage}:
if validation was recorded in the same club or with the same camera angles as training, the
model has \textbf{overfitted to one particular setup}.

To isolate the problem:

\begin{itemize}
  \item \textbf{Training vs.\ validation.} They should be \textbf{two genuinely independent
    subsets}. If they are not, validation loss tracks training loss very closely and we cannot
    tell when overfitting happens or whether the run is healthy.
  \item \textbf{Validation vs.\ test.} Look for what changes between them: \textbf{scale,
    perspective, FoV, resolution, color of the ball / field}, etc.
  \item \textbf{Probing.} A cheap check: take the features extracted by the trained model and
    train a \textbf{tiny classifier} to predict what I suspect is breaking it (the padel
    field, the camera resolution, the perspective). If that classifier is accurate, the
    representation \textbf{encodes this feature} and the decisions may be affected by it.
  \item \textbf{Annotate a few frames from the failing club.} The system is a pipeline, so
    the most direct test is to label a few hundred frames there and measure the two stages
    separately: first the pose estimation on its own, and then the hit detection fed with those
    ground-truth keypoints instead of the predicted ones. If the pose is already wrong, the
    problem is in the perception stage. If the pose is fine and the hits are still wrong, the
    problem is in what consumes it.
\end{itemize}

\section*{Q2. Impact of inconsistent annotations and label-quality strategy}

\subsection*{Impact on the model}

If annotators label differently, the model has \textbf{no way of knowing} whether the elbow
should go slightly higher or lower: it depends only on the annotator, and we give the model
no information about \textbf{who labelled what}. To minimise the loss it will predict what
the \textbf{average annotator} considers correct, but it can also pick up \textbf{spurious
correlations} that make predictions unreliable. The same happens with the
\textbf{timestamps}: label quality has a \textbf{direct impact} on performance, and with
random noise as labels the model would learn nothing useful. \textbf{Hit classification} is
affected too, since the hit type is probably inferred from wrist and elbow relative to the
racket at the moment of contact, and how much will depend on the \textbf{architectural
choices}.
Disagreement between annotators also sets a \textbf{noise floor}: the model cannot be more
consistent than the people who labelled the data, so measuring how much they disagree tells us
when it stops being worth improving the model and we should improve the labels instead.

\subsection*{How to detect it, and what to do}

Everything depends on whether we \textbf{know who annotated what} (for example, one folder
or file naming per annotator).

\begin{itemize}
  \item \textbf{If we know it}, we can look for systematic differences, e.g.\ the
    \textbf{average forearm length}. An annotator who places the wrist in the middle of the
    hand gives a consistent offset in every clip, so we can not only detect the bias but
    \textbf{correct it}, shortening the elbow-to-wrist vector by that amount. Same for timestamps:
    if one annotator always marks the hit with the ball at distance $x$ from the racket and
    another at distance $y$, that is again a correctable bias.
  \item \textbf{If we do not know it}, it gets more complicated. Then the decision depends on
    how \textbf{bad} the annotations are, how \textbf{expensive} relabelling would be, and how
    much performance drops if we simply \textbf{filter those parts out} (an \textbf{ablation
    run} makes sense here).
\end{itemize}

\textbf{Reweighting} adds complexity and possible \textbf{over-engineering} (you have to
choose a weight per partition), so I would not go there until everything else is exhausted.
I would also watch \textbf{sub-metrics for smashes} and check whether (1) overall precision
and (2) smash precision improve or get worse.

\section*{Q3. Handling rare but important smash events}

\subsection*{First, without retraining}

The cheapest thing is to experiment with the \textbf{decision threshold}: lowering it raises
recall on rare cases \textbf{with no retraining}, and the model may already know more about
smashes than the \textbf{argmax} suggests. The trade off is \textbf{more recall, less
precision} on smashes. If this is not enough, we retrain and modify the dataset or the loss.

\subsection*{The dataset}

If smashes are very rare, the model can reach a very good training loss \textbf{without ever
fitting them}. Augmenting them makes the model see these cases more often and \textbf{forces
it to fit this hit type}:

\begin{itemize}
  \item \textbf{Oversampling} (repeating clips): no new information, only more gradient
    weight, and it multiplies the risk of \textbf{memorising those clips}.
  \item \textbf{Augmentation}: mirroring, cropping / scaling, adding noise, etc.
  \item \textbf{Acquiring real smashes}: this is the \textbf{best solution}, since it is
    the only one adding \textbf{real new information} (and it lets me go play padel). A
    cheap way to get them from footage we already have is to run the current model with a
    \textbf{deliberately low smash threshold}, so it over-predicts smashes, and send only those
    candidate windows to be checked and annotated. The annotators then spend their time on
    smashes instead of on the footage that contains none. If we do not have any extra footage:
    book a padel court, practise smashes with friends, record them and add them to the dataset.
    I offer myself as a volunteer for the job :).
\end{itemize}

We can also update the \textbf{loss} so that a missed smash is \textbf{penalised more
heavily}.

\subsection*{Evaluation}

Here I would be \textbf{more careful}, because any change makes the results \textbf{not
comparable} with other versions of the model.

\begin{itemize}
  \item If validation already has \textbf{enough smashes} and they are correctly labelled,
    \textbf{leave it as it is}.
  \item If there are too few, the conclusions may not be \textbf{statistically significant}.
    In that case I would add smashes to validation as well, but only by \textbf{acquiring real
    ones},
    never by oversampling or augmentation. Loss results would not be comparable anymore.
  \item Keeping evaluation \textbf{untouched} is what lets me tell whether the training
    changes are real or just \textbf{overfitting artifacts}: the same augmentation on both
    sides could make me believe the model improved when it is only overfitting.
\end{itemize}

\subsection*{Trade-offs}

Weighting smashes too heavily will probably \textbf{degrade the other hit types}, so how much
to weight them depends on \textbf{how important they are for the product}. Augmentation costs
\textbf{training time} and changing the loss can add \textbf{instability}. If training time is
not a bottleneck, I would go with \textbf{augmenting the dataset}.

\section*{Q4. Stable-but-less-accurate vs.\ accurate-but-jittery model}

If two runs give noticeably different results, I would first train \textbf{more models with
different seeds}, check that they \textbf{converged properly} and understand what randomness
produced this.

For production I would \textbf{ship the stable one}: it is easier to justify a model with
lower precision than one that from time to time gives a completely \textbf{unpredictable
response}. It depends of course on how much precision we lose and how often the nonsense
happens, but for a product a \textbf{predictable model is better}, because you have to know
what you are selling and be able to speak for it when it makes mistakes.

For development I am \textbf{more interested in the jittery model}:

\begin{itemize}
  \item Which scenarios \textbf{break} it, and when exactly it produces inconsistent timing or
    jittery keypoints.
  \item Are those jitters \textbf{justified}? If they are (the ball is not visible, it is
    occluded, or similar), that would \textbf{change my decision} between the two models.
  \item How do we even measure the jitter? \textbf{Per-frame keypoint error cannot see it},
    because it compares every frame with its ground truth independently, so a model whose
    predictions oscillate between consecutive frames can still look better on average. To see it
    we have to compare each frame \textbf{with the previous timesteps}: displacement and jerk
    between consecutive frames, especially in segments where the player is almost still and the
    real movement should be close to zero.
  \item Can they be removed in \textbf{postprocessing}? A person moves smoothly and cannot be
    ``teleported'', so keypoints that jump can be identified as \textbf{noise} and filtered
    out.
\end{itemize}

\section*{Q5. Proposed model for this application}

I would split the architecture into \textbf{four different problems}:

\begin{enumerate}
  \item \textbf{Pose estimation of the players.} An \textbf{off-the-shelf model}
    (\href{https://arxiv.org/abs/2312.07526}{RTMO} or
    \href{https://huggingface.co/Ultralytics/YOLO26}{YOLO26-pose}), evaluated to see how it
    performs and whether \textbf{fine-tuning} is needed,
    plus lightweight \textbf{multi-person tracking} so keypoints carry a stable identity.
    I would lean towards RTMO because it is \textbf{one-stage}: it finds all the
    keypoints in the image in a single forward pass, so its
    cost does not depend on how many people are in the frame. A classic \textbf{top-down}
    model works the other way around, because it first detects each person and then runs a
    separate pose network on every crop. With four players on a padel court that means running
    the pose network four times per frame.
  \item \textbf{Ball tracking.} It probably needs a \textbf{higher frequency} than the pose
    estimation, or a couple of frames at a time, since the ball is a \textbf{tiny object}
    moving much faster. I would take a model \textbf{pretrained on tennis} or a similar sport
    and fine-tune it for padel.
  \item \textbf{Racket tracking.} Same approach, pretrained plus fine-tuning. I am \textbf{not
    sure it is necessary}, so I would first evaluate whether the racket can be inferred
    accurately enough from the \textbf{hand position}.
  \item \textbf{Hit classifier.} Trained on the \textbf{outputs of the previous stages}
    (poses, ball and racket positions), a \textbf{tiny ML model} is enough to determine the
    hit type, and it can be \textbf{image agnostic}.
\end{enumerate}

For 1, 2 and 3 I would pick models that run \textbf{fast on CPU}, ideally with no GPU at all,
which is why I would prefer a \textbf{CNN backbone} over a transformer one: vision
transformers are \textbf{less CPU friendly}. The main issue would probably be
\textbf{licensing}: checking which models we are allowed to use.

\end{document}
